%% ICLR 2027 submission skeleton — chemical-transformer expansion.
%%
%% POLICY CHECKLIST (executed by hand, not by LaTeX):
%%  [ ] Sep 16: freeze conservative abstract (below), draft on file.
%%  [ ] Sep 17: register abstract on OpenReview. Abstracts stay editable
%%      until Sep 25 — sharpen [X]-marked claims as P3/P4 land.
%%  [ ] Verify OpenReview profile accuracy (desk-reject grounds if wrong).
%%      Author list frozen at Sep 18; no adds. One-paper cap: this is it.
%%  [ ] 6-10 pages main text HARD (desk reject at 11). Appendix unlimited.
%%  [ ] Swap iclr2026_conference.sty -> official iclr2027 when published.
%%  [ ] RAAAI version is non-archival: no dual-submission conflict.
%%
%% PLAN A/B INTRO TOGGLE: decide ~Sep 12 on P0/P1 evidence.
%%   \introPlanBtrue  -> gate-gradient leads (P0 hit)
%%   \introPlanBfalse -> failure-diagnosis leads (workshop framing)
\newif\ifintroPlanB
\introPlanBfalse
\newcommand{\XX}{\textbf{\color{red}??}}          % pending number
\newcommand{\todoX}[1]{\textcolor{red}{[X: #1]}}  % pending sharpening

\documentclass{article}
\usepackage{iclr2027_conference}
\usepackage{tikz}
\usepackage{amsmath,amssymb}
\usepackage{xcolor}
\usepackage{booktabs}
\usepackage{graphicx}
\usepackage{multirow}
% auto-generated by paper/generate_tables_iclr.py — do not edit
\expandafter\def\csname GategradAcc\endcsname{74.5$\pm$1.0}
\expandafter\def\csname ShuffledAcc\endcsname{74.7$\pm$1.1}
\expandafter\def\csname TagAcc\endcsname{74.8$\pm$1.3}
\newcommand{\pzerotable}[0]{
\begin{tabular}{lcccc}
\toprule
Stage-2 arm & Acc & Balanced & MFLOPs/tok & Gate spread \\
\midrule
Tag targets (labels) & 74.8$\pm$1.3 & 83.6$\pm$0.8 & 19.4$\pm$0.1 & 0.21$\pm$0.06 \\
Gate-grad (measured value) & 74.5$\pm$1.0 & 83.4$\pm$0.7 & 19.2$\pm$0.0 & 0.08$\pm$0.01 \\
Shuffled control & 74.7$\pm$1.1 & 83.5$\pm$0.7 & 19.3$\pm$0.0 & 0.06$\pm$0.01 \\
\bottomrule
\end{tabular}
}
\newcommand{\budgettable}[0]{
\begin{tabular}{lcccc}
\toprule
Arm & mean target & nominal spread & realized $\bar g$ & realized spread \\
\midrule
Tag targets & 0.51 (fixed) & 0.60 & 0.51 & 0.21$\pm$0.06 \\
Gate-grad & ? & 0.70 & 0.51 & 0.08$\pm$0.01 \\
Shuffled & 0.30 & 0.70 & 0.51 & 0.06$\pm$0.01 \\
\bottomrule
\end{tabular}
}
\newcommand{\migrationtable}[0]{
% migration table pending
}
\newcommand{\windowtable}[0]{
% window table pending
}
   % auto-generated macros + table bodies

\iclrfinalcopy % submissions compile with this; camera-ready identical

\title{Endogenous Signals Fail to Allocate: \\\ Measured Signals Supervise
Token-Level Adaptive Compute at Scale}

\author{Anonymous authors\\
Paper under double-blind review}

\newcommand{\LN}{\mathrm{LN}}
\newcommand{\FFN}{\mathrm{FFN}}
\newcommand{\attn}{\mathrm{attn}}

\begin{document}

\maketitle

\begin{abstract}
Token-level adaptive compute usually keys off endogenous signals---attention
entropy, router scores, confidence---assumed to track per-token difficulty. On
arithmetic with graded difficulty we show this assumption fails in three
distinct ways: an entropy-driven gate collapses to a constant, a tag-initialized
gate inverts (allocating \emph{more} compute to easy tokens), and an
unsupervised learned head is vacuous. Explicit per-token difficulty supervision
with a budget constraint fixes all three: correct allocation ordering in every
seed ($r = 0.95$) and full-compute accuracy at ${\sim}42\%$ fewer FLOPs per
token. A pre-registered shuffled-targets
control then separates two explanations: at testbed scale, supervised
\emph{variation} under an explicit budget---not target content---carries
the benefit. This makes the content question empirical, and we pursue it
with a direct measurement: the marginal value of compute for every token,
$\partial\mathcal{L}/\partial g_t$, obtainable in one backward pass. We
track how measured value distributes across difficulty bins as models scale
\todoX{42M $\to$ 400M}, and impose measured targets on natural-language
domains (web text, code, math) where difficulty must be measured rather
than given, on pretrained \todoX{1.4B--2.8B} models under
iso-training-FLOPs controls.
\end{abstract}

\section{Introduction}
\label{sec:intro}

\ifintroPlanB
% ---- PLAN B: gate-gradient leads (use if P0 hits across seeds) ----
Adaptive-compute transformers promise to spend FLOPs where they help. Every
practical method must answer two questions per token: \emph{how much compute
does this token deserve?} and \emph{what signal decides?} The field's default
answer to the second question is an endogenous one---attention entropy, router
logits, confidence---trained end-to-end and trusted to discover difficulty.
Our first result is that this trust is misplaced: driven end-to-end, every
endogenous signal we test fails to allocate in a characteristic way (collapse,
inversion, or vacuity; \S\ref{sec:failures}).
Our second result is constructive and, we argue, the more important one. The
object these signals are groping for is not difficulty but the \emph{value of
compute}: the marginal loss reduction one more unit of compute would buy. It
does not need to be discovered by training. For a trained model it can be
\emph{measured}: one backward pass with the per-token FFN gates treated as
differentiable leaves yields $\partial \mathcal{L}/\partial g_t$ for every
token and layer at once. \emph{Compute-value supervision}---imposing this
measured signal on the gate under an explicit budget---allocates as well as
ground-truth difficulty labels, and better than the strongest
static-schedule control, in every setting we test---including
natural-language domains where no difficulty labels exist at all. The two
supervisions are not the same object, and the difference is measurable: loss
says ``hard,'' the gate gradient says ``helpable.'' Difficulty was always a
proxy for value; end-to-end training chases the proxy uncontrollably (it
overfeeds easy tokens because that is where loss is most reducible), while
the budget term makes chasing value safe. The practical consequence is a
recipe: train dense, measure per-token compute value in one backward pass,
impose allocation offline---one extra backward pass, no search for a magic
endogenous signal.
% ---- end Plan B ----
\else
% ---- PLAN A: failure diagnosis leads (workshop framing) ----
Token-level adaptive compute lets a transformer spend FLOPs where they help.
Methods that decide \emph{which} tokens deserve compute typically key off
endogenous signals---attention entropy, router logits, confidence---whose
correlation with task difficulty is assumed rather than tested. We test the
assumption on arithmetic with graded difficulty and per-token ground truth,
and find it false in three distinct ways: entropy-driven gates collapse to a
constant; tag-initialized gates \emph{invert}, allocating more compute to easy
than hard tokens; unsupervised learned heads stay vacuous. Explicit
difficulty supervision with a budget constraint fixes all three.
We then ask what the fix really requires, and the answer sharpens the
concept: the useful object is not difficulty but the \emph{value of compute}
(the marginal loss reduction one more unit of compute would buy). Difficulty
labels are one proxy for it; the trained model itself supplies a better one.
Where labels are absent we substitute \emph{measured} signals computed
offline---per-token loss, and $\partial\mathcal{L}/\partial g_t$, obtained in
one backward pass with gates as differentiable leaves. Compute-value
supervision allocates as well as difficulty labels wherever we test, while
the same signals trained end-to-end do not: the failure is in the training
dynamics, not the information source.
% ---- end Plan A ----
\fi

\paragraph{Contributions.}
\begin{enumerate}
\item \textbf{A three-way failure diagnosis} of endogenous allocation signals
(collapse, inversion, vacuity), each with a mechanistic account
(\S\ref{sec:failures}); a standard Mixture-of-Depths router fails
analogously.
\item \textbf{Compute-value supervision}: targets derived offline from the
trained model itself---per-token loss, and $\partial\mathcal{L}/\partial g$
in one backward pass---match label-supervised allocation, with a
shuffled-targets control (and a pre-registered decision rule) separating
the specific measured targets from supervised variation alone
(\S\ref{sec:measure}, \S\ref{sec:p0}); per-bin value signatures migrate
with model scale as the value-not-difficulty distinction predicts
(\S\ref{sec:scale}).
\item \textbf{Iso-FLOPs methodology}: all method comparisons at matched
\emph{billed training FLOPs} (gated billing, forward + backward), with
within-segment interpolation and a budget sweep; the workshop version's
``similar total budget'' caveat is retired (\S\ref{sec:isoflops}).
\item \textbf{Scale and transfer}: the effect survives scaling from 11M to
\todoX{1B} parameters on arithmetic and transfers to natural language
(web text, code, math) on pretrained \todoX{1.4B--2.8B} models where
difficulty is measured, not given (\S\ref{sec:scale}, \S\ref{sec:nl}).
\item \textbf{Realization}: a tiered-FFN serving benchmark showing
tokens/sec scaling with the mean gate, converting billed-FLOPs savings into
wall-clock (\S\ref{sec:realization}).
\end{enumerate}

\section{Related Work}
\label{sec:related}

\paragraph{Adaptive compute.} Adaptive Computation Time
\citep{graves2016act}, early-exit and cascades, and Mixture-of-Depths
\citep{raposo2024mixture} route tokens by learned signals; MoE+router work
attaches learned routers to expert selection. \citet{raposo2024mixture} train
routers from scratch; our MoD baseline on the language-model track inherits
the resulting handicap---a router trained during a short fine-tune, not a
full pretraining run---which we acknowledge as a baseline limitation, not a
strawman.
\paragraph{Conditional computation and mixture-of-experts.}
\paragraph{Entropy and confidence as uncertainty proxies.} Draft-model
disagreement in speculative decoding is a deployed, per-token difficulty
signal: acceptance probability measures ``hard for a small model,'' and our
gate-gradient targets are the same quantity computed for layers instead of
models.
\paragraph{Resource-aware agents.}

\section{Method}
\label{sec:method}

\subsection{Architecture and per-token gating}
\label{sec:arch}

Each layer computes
\begin{equation}
x_2 = \LN\!\bigl(x_1 + \attn(x_1)\bigr), \qquad
x_3 = \LN\!\bigl(x_2 + g_t \cdot \FFN(x_2)\bigr),
\label{eq:gate}
\end{equation}
where $g_t \in [0,1]$ is a per-token, per-layer gate produced by a small
per-layer \emph{DifficultyHead} (a two-layer MLP on the post-attention state
$x_2$, sigmoid-squashed). Compute is billed proportionally: a position gated
at $g$ is billed $g$ times the per-token FFN cost; attention is always paid
in full. All variants share Eq.~\eqref{eq:gate} and the accounting, so
comparisons isolate the gating signal. We gate the FFN only, never the
attention's keys or values: skipping a token's FFN is safe because downstream
tokens still attend to it, whereas skipping K/V irreversibly deletes
information that later positions may need.

\subsection{Why endogenous signals fail: three diagnostics}
\label{sec:failures}

End-to-end training offers the gate one channel of feedback: the task loss.
\emph{Collapse}: a signal with no difficulty information (attention entropy,
$r \approx -0.06$ per token with difficulty) drives the gate to a constant.
\emph{Inversion}: a ground-truth initialization makes the gate vary, but easy
problems are memorizable, so their cross-entropy gradient is largest and
training allocates them \emph{more} compute (per-problem gate--difficulty
$r = -0.59$). \emph{Vacuity}: a head trained only through the task loss
receives no signal about what difficulty is; the gate stays flat and unstable
across seeds. We report the two iteration-1 quantities as distinct: per-token
entropy--difficulty correlation $r \approx -0.06$, and per-problem
entropy-gate--difficulty $r = 0.23$ (the gate passes the tag embedding through
a bottleneck; the two correlations measure different objects).

\subsection{Explicit difficulty supervision with a budget constraint}
\label{sec:supervised}

Training is two-stage. \textbf{Stage 1} trains the body conventionally with
the gate fixed at $1$ and checkpoints it. \textbf{Stage 2} attaches the
DifficultyHeads and minimizes
\begin{equation}
\mathcal{L} = \mathcal{L}_{\mathrm{CE}}
+ \lambda_{\mathrm{mse}}\,\tfrac{1}{N}\sum_{t}\bigl(\sigma(s_t) - \tau_t\bigr)^2
+ \lambda_{\mathrm{budget}}\,\bigl(\bar{g} - 0.5\bigr)^2,
\label{eq:loss}
\end{equation}
with $\lambda_{\mathrm{mse}} = 5.0$, $\lambda_{\mathrm{budget}} = 0.5$; the
body trains at a near-frozen learning rate ($10^{-2}\times$ the heads').
Stage 1 is shared by every fork of the method (all Stage-2 arms load the same
checkpoint per seed), so method comparisons are paired by body.

\subsection{Measuring the value of compute}
\label{sec:measure}

Where difficulty labels are absent (all natural-language settings), targets
$\tau_t$ are \emph{measured offline} from the Stage-1 checkpoint:
\begin{equation}
\mathrm{score}_t \;=\; -\sum_{\ell} \frac{\partial \mathcal{L}}{\partial
g_{t,\ell}}\Big|_{g=1},
\label{eq:gategrad}
\end{equation}
the marginal loss reduction from one unit of FFN compute at token $t$,
obtained for all tokens and layers in a single forward--backward pass with
gates as differentiable leaves. Scores ${\le}0$ are \emph{irreducible}
(compute does not help there) and receive the floor target $0.1$; positive
scores are quartile-binned to $0.2/0.4/0.6/0.8$. This is the paper's key
conceptual move: tags label \emph{difficulty}; Eq.~\eqref{eq:gategrad}
measures the \emph{value} of compute. The two are different objects, and the
difference is testable: score--loss correlation is near zero in our runs---%
loss says ``hard,'' the gate gradient says ``helpable''---and a token's tag
bin does not constrain its score (a hard bin mixes irreducible tokens with
worth-treating ones). We therefore expect, and report, that compute-value
supervision does \emph{not} order gates by difficulty bins at the testbed
scale: bin-level spread is not its target. The workshop-era inversion---%
end-to-end training overfeeding easy tokens---is the same phenomenon seen
from the other side: training chases reducibility wherever it lives, but
uncontrollably; the budget term is what makes chasing value safe. A
loss-quantile variant (the proxy-difficulty ablation) bins tokens by
per-token loss; a \emph{shuffled} control permutes targets within each
sequence, testing whether any varying supervision would do.

\paragraph{First-order oracle allocation.} The same measurements give a
reference point: sorting per-(token, layer) scores and filling until the
budget is exhausted yields the first-order-optimal allocation under a
\emph{frozen} body, evaluated zero-shot by overriding the gates at
inference. We label it accordingly: it bounds what allocation alone buys
without retraining. Trained methods may exceed it---their bodies adapt to
their gates---and we interpret that as adaptation compounding allocation,
not as a paradox.

\subsection{Iso-FLOPs training protocol}
\label{sec:isoflops}

Every run logs cumulative \emph{billed training FLOPs}: the gated billing of
\S\ref{sec:arch} applied to forward, with backward billed at twice forward
(total $3\times$). Methods are compared at matched achieved budgets by
linear interpolation of the evaluation history \emph{within} training
segments only---two-stage runs have a kink at the stage boundary, and budgets
common to all methods are capped at the shortest run's endpoint (no
extrapolation). At matched billed FLOPs, lower-mean-gate methods receive more
optimizer steps per billed dollar; that is the deployment claim under test,
and steps-matched comparisons are reported alongside in
Appendix~\ref{app:stepsmatched}.

\section{Experimental setup}
\label{sec:setup}

\paragraph{Arithmetic.} Two regimes with graded difficulty: \emph{tagged}
([E]/[M]/[H] prefixes; easy 1--2-digit $\pm$, medium 4--5-digit $\pm$, hard
$4{\times}4$, $5{\times}4$ and $6{\times}5$-digit multiplication at scale) and
\emph{mixed} (no tags; operands 1--7 digits; bin = larger operand's digit
count). The hard regime is chosen per scale so the largest model sits inside
a difficulty window: hard-bin accuracy neither saturated nor collapsed
(maintained in $[40\%, 85\%]$ by the calibration probes). Frozen held-out
sets, identical across methods and seeds; accuracy on answer tokens only.
Models: an 11M testbed (the workshop setting, exactly reproduced) and
\todoX{150M/400M/1B} scaled models (RoPE, flash attention, bf16, cosine
schedule; identical across methods).

\paragraph{Natural language.} Pretrained \todoX{Pythia-1.4b / 2.8b}
continue-trained on three domains where per-token difficulty is latent:
web text (FineWeb-Edu), Python code (CodeSearchNet), and chain-of-thought
math (MetaMathQA). Stage-1 dense continue-training is shared by all forks;
the oracle of \S\ref{sec:measure} runs once over a pool of training
sequences plus the held-out split, producing per-token grad bins, loss bins,
and diagnostics. Metrics: held-out loss at matched billed inference
FLOPs/token, per-bin loss and teacher-forced accuracy, gate--bin ordering
$r$, and gate spread. We report pass@k only as texture: a fine-tuned
\todoX{1.4B} model scores single digits on HumanEval-style evaluation, and
per-bin teacher-forced metrics are the primary outcome.

\paragraph{Methods compared.} \emph{Baseline} (dense; = Stage-1), ours
(\emph{gate-grad} targets; \emph{loss-quantile} ablation), \emph{shuffled}
control, \emph{Fixed-schedule} (static $g{=}0.5$), \emph{MoD} at matched
capacity, \emph{Chemical-off} (architecture control), and the endogenous
iterations (\emph{entropy}, \emph{random}, \emph{tag-only}) at testbed scale.

\section{Results}
\label{sec:results}

\subsection{Endogenous signals fail; supervision fixes allocation}
\label{sec:failures-results}

Table~\ref{tab:testbed} and Figure~\ref{fig:journey}: \XX. (Ported from the
workshop version: entropy $0.232/0.233/0.233$ across E/M/H; tag-init
inverted; unsupervised vacuous; supervised spread $+0.20$, ordering
$r = 0.95$ in every seed.)

\subsection{Measured value matches labels at the testbed}
\label{sec:p0}

\begin{table}[t]
\centering
\caption{P0: paired Stage-2 arms from the same Stage-1 body per seed
(3 seeds, mean$\pm$std). Billed training FLOPs are identical across the
bottom two arms by construction (the oracle pass is billed to both); the
tag arm has no oracle pass.}
\label{tab:p0}
\pzerotable
\end{table}

\paragraph{Pre-registered decision rule.} Written before any shuffled number
existed: the measured-targets claim survives iff the gate-grad arm beats the
shuffled control by ${\ge}1.0$ point of held-out accuracy, paired by seed,
at matched billed training FLOPs, across all three seeds. No post-hoc
reinterpretation either way.

\paragraph{Outcome: the rule fired against the measured targets.} Paired
gate-grad $-$ shuffled accuracy across seeds: $-0.003$, $-0.001$, $-0.000$
(Table~\ref{tab:p0}).
At testbed scale, \emph{any} varying supervision under the budget
constraint---even targets permuted into noise---recovers the same accuracy
($0.747 \pm 0.009$ shuffled vs.\ $0.745 \pm 0.008$ gate-grad vs.\ $0.748 \pm
0.010$ tag) and the same flat gate structure, while all three dominate a
constant gate trained from scratch by ${\sim}2$ points at the same inference
budget. Two honest readings follow. First, the workshop result sharpens: at
$11$M, what supervision supplies is not allocation quality but a dense-then-
gate protocol whose accuracy is insensitive to target content; allocation
structure (the tag arm's $+0.21$ spread) is real but not
accuracy-load-bearing here. Second, the value-of-compute concept is not
refuted---it is falsifiable at scale: if Eq.~\eqref{eq:gategrad} measures
reducible value rather than noise, its per-bin signature must migrate as
model capability grows (\S\ref{sec:scale}), and under it the natural-language
targets differ from shuffled ones precisely where compute has leverage
(\S\ref{sec:nl}). Both tests are pre-registered; \XX\ reports their outcome.

\begin{table}[t]
\centering
\caption{Supervision targets vs.\ realized gates (budget redistribution).
The floor + quartile targets have mean $\approx 0.31$ and nominal spread
$0.7$; the pinned budget ($\bar g \to 0.5$) lifts and compresses them. The
compression is the budget term redistributing, not the head failing to
allocate --- a pinned mean is what makes the inference budget a settable
dial.}
\label{tab:budget}
\budgettable
\end{table}

Bin-level spread is \emph{not} the allocation metric for measured targets
(\S\ref{sec:measure}): gate-grad varies within problems (a hard bin contains
both irreducible and worth-treating tokens), so its E/M/H spread is smaller
while its per-token allocation is finer; we report gate--score agreement on
held-out tokens (\XX) and per-token rather than per-bin summaries.

\subsection{Scale: the effect survives 11M $\to$ \todoX{1B}}
\label{sec:scale}

\paragraph{Pre-registered readout: value migrates with capability.} If the
oracle measures value rather than difficulty, its per-bin signature must
change with model scale: an 11M model finds most hard-bin tokens
irreducible, but a $400$M model does not. Prediction (written before the P1
probes ran): per-bin mean gate-grad score on the hard bin rises, and its
fraction of non-positive scores falls toward the medium bin's, across
$42$M $\to$ $150$M $\to$ $400$M. A flat signature at \$10 of probe cost
would have falsified the irreducibility mechanism before the expensive
stage. Result: Figure~\ref{fig:value-migration} --- \XX.

Figure~\ref{fig:frontier-scale}: accuracy vs.\ billed inference FLOPs at
matched training budgets across \todoX{11M/150M/400M/1B}. \XX

\subsection{Natural language: measured difficulty transfers}
\label{sec:nl}

Table~\ref{tab:nl}: three domains on \todoX{Pythia-1.4b}. \XX
Diagnostics: fraction of irreducible tokens (score ${\le}0$) \XX;
score--loss correlation \XX\ (measured value is not difficulty: high-loss
tokens can be hopeless); bin--position correlation \XX\ (early tokens
mechanically influence more downstream loss; reported, and the doc-boundary
tokens landing in the floor bin is the expected irreducibility signature);
token-type breakdown of the top bin \XX.

\subsection{Where does the compute go? Oracle gap and layer profile}
\label{sec:oracle-gap}

Against the first-order frozen-body oracle (\S\ref{sec:measure}): \XX.
Per-layer gate profiles by difficulty bin: \XX.

\subsection{Wall-clock realization}
\label{sec:realization}

Billed FLOPs become real savings only under sparse execution. A tiered-FFN
serving benchmark (tokens sorted by gate; active-set gathered GEMM) shows
tokens/sec scaling as \XX\ at mean gate $0.5$ on the \todoX{1B} shape.

\section{Discussion}
\label{sec:discussion}

\paragraph{Implications.} The recipe---train dense, measure, impose---needs
no endogenous-signal search and plugs into serving stacks that already
produce per-token difficulty estimates (speculative-decoding acceptance).
\paragraph{When static schedules suffice.} \XX
\paragraph{Limitations.} Arithmetic difficulty is operand-size-defined at
testbed scale; the natural-language bins inherit the reference model's
notion of difficulty. The two-stage body is near-frozen by default; a
full-unfreeze ablation is reported in Appendix~\ref{app:ablations}. Billed
FLOPs are an accounting device; \S\ref{sec:realization} is the bridge to
wall-clock, and realized speedups depend on kernel support. Gate-target
quartiles are hand-set; Eq.~\eqref{eq:gategrad} replaces their content but
not their form.

\section*{Reproducibility statement}
All code, manifests, and per-run JSONs (config, history, billed-FLOPs
curves) are in the anonymized repository; every table and figure regenerates
from committed artifacts. Seeds, held-out sets, and the FLOPs accounting are
specified in \S\ref{sec:setup} and Appendix~\ref{app:details}.

\section*{Compute statement}
Total rented-GPU spend: \todoX{\$\XX} (\todoX{\XX} A100-hours), itemized per
stage in Appendix~\ref{app:compute}; testbed runs on an Apple M3 (MPS).

\bibliography{iclr2027_conference}
\bibliographystyle{iclr_conference}

\appendix

\section{Steps-matched comparisons}
\label{app:stepsmatched}
\XX

\section{Ablations}
\label{app:ablations}
Full-unfreeze Stage-2; per-layer target freedom; shuffled control at all
scales. \XX

\section{Experimental details}
\label{app:details}
\XX

\section{Compute ledger}
\label{app:compute}
\XX

\end{document}
